% =====================================================================
%  FICHE 8 — THÈME 3 — NIVEAU DIFFICILE — ÉNONCÉS
% =====================================================================
\documentclass[11pt,a4paper]{article}
\usepackage[utf8]{inputenc}
\usepackage[T1]{fontenc}
\usepackage{lmodern}
\usepackage[french]{babel}
\usepackage{amsmath,amssymb,mathtools}
\usepackage{icomma}
\usepackage{siunitx}
\sisetup{output-decimal-marker={,},per-mode=power,group-separator={\,},group-minimum-digits=5}
\usepackage[version=4]{mhchem}
\usepackage{xcolor}
\usepackage{tikz}
\usepackage[most]{tcolorbox}
\usepackage{enumitem}
\usepackage{array}
\usepackage{fancyhdr}
\usepackage{caption}
\captionsetup{labelformat=empty,justification=centering,font=small}
\usepackage[hidelinks]{hyperref}
\usetikzlibrary{arrows.meta,positioning,calc}
\usepackage[a4paper,margin=2.1cm,headheight=26pt,headsep=10pt]{geometry}
\usepackage{titlesec}

\definecolor{primary}{RGB}{150,98,162}
\definecolor{primarydark}{RGB}{92,52,110}
\definecolor{chaleur}{RGB}{214,72,33}
\definecolor{endocol}{RGB}{0,110,180}

\tcbset{boxbase/.style={enhanced,breakable,boxrule=0.9pt,arc=2.5pt,left=3mm,right=3mm,top=2.3mm,bottom=2.3mm,fonttitle=\bfseries,coltitle=white,toptitle=0.6mm,bottomtitle=0.6mm}}
\newtcolorbox[auto counter]{exercice}[1][]{boxbase,colback=primary!4,colframe=primary,title={\textsc{Exercice}~\thetcbcounter\ \ifx\relax#1\relax\relax\else\textmd{--- \itshape #1}\fi}}

\pagestyle{fancy}
\fancyhf{}
\fancyhead[L]{\small\textcolor{primarydark}{\textbf{Fiche 8} — Thème 3 : Calorimétrie}}
\fancyhead[R]{\small\textcolor{primarydark}{\textsc{Difficile}}}
\fancyfoot[C]{\small\thepage}
\renewcommand{\headrulewidth}{0.8pt}
\renewcommand{\headrule}{\hbox to\headwidth{\color{primary}\leaders\hrule height \headrulewidth\hfill}}
\setlength{\parindent}{0pt}\setlength{\parskip}{3pt}

\begin{document}

\begin{center}
\begin{tikzpicture}
\node[rectangle,rounded corners=7pt,fill=primary,text=white,minimum width=16cm,
      minimum height=1.1cm,align=center,inner sep=6pt]
  {{\large\bfseries Thème 3 — Calorimétrie}\\[1pt]{\small Niveau \textsc{Difficile} — énoncés}};
\end{tikzpicture}
\end{center}

\begin{exercice}[de la glace à la vapeur — bilan complet]
On veut transformer $m=\SI{50}{\gram}$ de glace prise à \SI{-40}{\celsius} en vapeur d'eau portée à
\SI{120}{\celsius}. Données : $c_{\text{glace}}=\SI{2.060}{\kilo\joule\per\kelvin\per\kilogram}$,
$c_{\text{eau}}=\SI{4.185}{\kilo\joule\per\kelvin\per\kilogram}$,
$c_{\text{vap}}=\SI{1.85}{\kilo\joule\per\kelvin\per\kilogram}$,
$L_{\text{fus}}=\SI{334}{\kilo\joule\per\kilogram}$, $L_{\text{vap}}=\SI{2257}{\kilo\joule\per\kilogram}$.
\begin{center}
\begin{tikzpicture}[scale=0.82,every node/.style={font=\scriptsize}]
  \draw[->,line width=1pt,endocol] (0,0) -- (2.0,0); \node[above] at (1.0,0){\textcircled{1}};
  \node[below left] at (0,0){$-40\,\si{\celsius}$};
  \draw[->,line width=1pt,chaleur] (2.0,0) -- (2.0,0.9); \node[right] at (2.0,0.45){\textcircled{2}}; \node[below] at (2.0,0){$0\,\si{\celsius}$};
  \draw[->,line width=1pt,endocol] (2.0,0.9) -- (4.6,0.9); \node[above] at (3.3,0.9){\textcircled{3}};
  \draw[->,line width=1pt,chaleur] (4.6,0.9) -- (4.6,1.8); \node[right] at (4.6,1.35){\textcircled{4}}; \node[below] at (4.6,0.9){$100\,\si{\celsius}$};
  \draw[->,line width=1pt,endocol] (4.6,1.8) -- (6.7,1.8); \node[above] at (5.65,1.8){\textcircled{5}};
  \node[above right] at (6.7,1.8){$120\,\si{\celsius}$};
\end{tikzpicture}
\end{center}
\begin{enumerate}[label=\alph*),leftmargin=2em,topsep=2pt,itemsep=2pt]
  \item Nommer les cinq étapes \textcircled{1}--\textcircled{5} et indiquer la formule utilisée pour chacune.
  \item Calculer la chaleur de chaque étape.
  \item En déduire la chaleur totale $Q$.
  \item Quelle étape domine le bilan ? Quelle part (en \%) du total représente-t-elle ?
\end{enumerate}
\end{exercice}

\begin{exercice}[barre de fer froide dans l'eau — congélation partielle]
Un calorimètre adiabatique (capacité négligeable) contient $m_{\text{eau}}=\SI{200}{\gram}$ d'eau à
\SI{5}{\celsius}. On y plonge une barre de fer de $m_{\text{fer}}=\SI{500}{\gram}$ sortant d'un
congélateur à \SI{-34}{\celsius}. Données : $c_{\text{eau}}=\SI{4}{\kilo\joule\per\kelvin\per\kilogram}$,
$c_{\text{fer}}=\SI{0.5}{\kilo\joule\per\kelvin\per\kilogram}$, $L_{\text{fus}}=\SI{334}{\kilo\joule\per\kilogram}$.
\begin{enumerate}[label=\alph*),leftmargin=2em,topsep=2pt,itemsep=2pt]
  \item Quelle énergie faut-il fournir au fer pour l'amener de \SI{-34}{\celsius} à \SI{0}{\celsius} ?
  \item Quelle énergie l'eau peut-elle céder en se refroidissant de \SI{5}{\celsius} à \SI{0}{\celsius} ?
  \item Comparer : montrer qu'une partie de l'eau doit \textbf{geler} pour boucler le bilan. Quelle
        énergie manque-t-il ?
  \item Vérifier que l'eau ne gèle pas entièrement, puis calculer la masse d'eau qui gèle.
  \item Donner la température finale et la masse d'eau restée liquide.
\end{enumerate}
\end{exercice}

\begin{exercice}[dissolution avec un calorimètre qui s'échauffe]
Dans un calorimètre en verre de masse $m_v=\SI{200}{\gram}$ contenant $m_e=\SI{150}{\gram}$ d'eau à
\SI{21.0}{\celsius}, on dissout $n=\SI{0.10}{\mol}$ de sulfate de cuivre(II) solide. La température monte
à \SI{23.5}{\celsius}. On admet que seule \textbf{la moitié} de la masse du verre s'échauffe. Données :
$c_e=\SI{4.185}{\kilo\joule\per\kelvin\per\kilogram}$, $c_v=\SI{0.84}{\kilo\joule\per\kelvin\per\kilogram}$.
\begin{enumerate}[label=\alph*),leftmargin=2em,topsep=2pt,itemsep=2pt]
  \item Calculer $\Delta T$.
  \item Calculer la chaleur $Q_1$ captée par l'eau.
  \item Calculer la chaleur $Q_2$ captée par le verre (attention : seule la moitié s'échauffe).
  \item En déduire la chaleur totale captée, puis l'enthalpie \textbf{molaire} de dissolution $\Delta H$
        (avec son signe).
\end{enumerate}
\end{exercice}

\end{document}
